% Final project proposal, ANLP -- ACL style.
%
% Build (Overleaf: pdfLaTeX + BibTeX, or just press Recompile):
%   pdflatex main && bibtex main && pdflatex main && pdflatex main
%
% Team: Claude's Plan. Rename the output PDF to ClaudesPlan-Proposal.pdf
%
% Requires, in the same directory: acl.sty, acl_natbib.bst, references.bib
%
\documentclass[11pt]{article}

% "preprint" = non-anonymous, with page numbers.
% Use "final" for a non-anonymous version without page numbers,
% or "review" for the anonymised, line-numbered version.
\usepackage[preprint]{acl}

\usepackage{times}
\usepackage{latexsym}
\usepackage[T1]{fontenc}
\usepackage[utf8]{inputenc}
\usepackage{microtype}
\usepackage{amsmath}
\usepackage{booktabs}
\usepackage{graphicx}

% acl.sty already loads natbib, geometry, caption and hyperref, and already
% sets \bibliographystyle{acl_natbib}. Do not load or set them again.

% The default title box (11 baselines) is too short for four authors plus
% roll numbers, team name, course, affiliation and the e-mail line.
\setlength\titlebox{6.4cm}

\title{Event-Aware Representation Learning for\\
Neutral Multi-Source News Summarization}

\author{
  \begin{tabular}{@{}cccc@{}}
    Chaitanya Thogata & Harshit Lalwani & Manas Agrawal & Raunak Seksaria \\
    {\normalfont (2023113018)} & {\normalfont (2023111028)} &
    {\normalfont (2023113023)} & {\normalfont (2023113019)} \\
  \end{tabular} \\[3pt]
  Team Name: Claude's Plan \\[1pt]
  {\normalfont Advanced NLP} \\[1pt]
  {\normalfont International Institute of Information Technology, Hyderabad} \\[1pt]
  {\normalfont\small\texttt{\{chaitanya.thogata, harshit.lalwani,}} \\
  {\normalfont\small\texttt{manas.agrawal, raunak.seksaria\}@research.iiit.ac.in}}
}

\begin{document}
\maketitle

\begin{abstract}
Recent work on neutral news generation suggests that shaping latent
representations may help mitigate framing bias. We propose to investigate
whether event identity can serve as the organizing principle of the
representation space. Our primary approach will use event-aware triplet
learning to bring representations of reports covering the same event closer
while separating carefully chosen contrasting reports, jointly with a neutral
text generation objective. Cross-report supervision will be used only during
training, while inference will operate on a single article. As an exploratory
extension, we also plan to study JEPA-style cross-report prediction, where the
representation of one same-event report is used to predict another. We will
evaluate whether these objectives improve neutrality, event retention, and
robustness to framing variation.
\end{abstract}

\section{Project Idea}

We propose a neutral multi-source news summarization system whose internal
representation is explicitly organised around the \emph{event} being reported
rather than the \emph{publisher} reporting it. Given several articles about the
same story from outlets of differing political leaning, the model produces a
single neutral summary; alongside the usual generation objective, its encoder is
trained with a triplet loss that pulls together differently-framed reports of the
same event and pushes apart same-publisher reports of different events.

The deliverable is a summarization model with an event-aware
representation-learning objective, evaluated on how much event content it
retains against how much publisher-specific information remains recoverable from
its representations.

Before writing this proposal we profiled every corpus we intend to use and
reimplemented the evaluation metric the relevant literature reports. Three of the
findings from that audit changed the design, and we state them here because they
are the substance of our motivation rather than background detail.

\section{Motivation}

News reports are not neutral observations of events. Different publishers may
emphasize different facts, use different language, omit details, or frame the
same occurrence through different editorial or ideological perspectives. This
makes neutralization difficult: a model must preserve the underlying event
while reducing reliance on source-specific presentation. Existing work
suggests that this problem is not purely about output wording. NeutraSum
directly shapes latent representations toward neutrality, while event-relation
graph methods show that explicitly modeling event information can improve
framing-bias mitigation and content preservation. These findings motivate us
to explore whether the shared information across differently framed reports
can provide a useful learning signal. Rather than treating any single source
as an authoritative account, multiple reports of the same occurrence can be
viewed as complementary views whose common content may better reflect the
information that should be preserved during neutralization.

\section{Preliminary Analysis}

We have already built the data pipeline and run the CPU-side measurements; the
numbers quoted above come from it. Table~\ref{tab:corpora} summarises what each
corpus can support, and Table~\ref{tab:probe} gives our preliminary probe
results.

\begin{table}[t]
\centering\small
\setlength{\tabcolsep}{4pt}
\begin{tabular}{@{}lccc@{}}
\toprule
\textbf{Corpus} & \textbf{Event} & \textbf{Pub.} & \textbf{Neutral} \\
\midrule
NeuS (3{,}095 st.)    & yes & yes & yes \\
AllSides (37{,}554)   & no  & yes & no  \\
WCEP-10 (10{,}200)    & yes & no  & partial \\
WNC (181k)            & no  & no  & partial \\
\bottomrule
\end{tabular}
\caption{Which labels each corpus carries: event grouping, publisher
(\textbf{Pub.}) and a neutral target. Sizes are stories, articles, clusters and
sentence pairs respectively. An event-aware triplet objective needs all three
labels on one example; only NeuS supplies them together, and it is the smallest
of the four.}
\label{tab:corpora}
\end{table}

\begin{table}[t]
\centering\small
\setlength{\tabcolsep}{4pt}
\begin{tabular}{@{}lrrr@{}}
\toprule
\textbf{TF-IDF probe} & \textbf{Acc.} & \textbf{Base} & \textbf{Cls.} \\
\midrule
Publisher, random split  & 17.8 & 9.2  & 229 \\
Stance, random split     & 40.2 & 37.3 & 3 \\
Stance, outlet-disjoint  & 35.4 & 46.0 & 3 \\
\bottomrule
\end{tabular}
\caption{Lexical probes on 9{,}282 NeuS articles; \textbf{Base} is the
test-majority baseline throughout and \textbf{Cls.} the number of classes.
Publisher is recoverable from word counts at roughly twice its baseline. Stance
is barely recoverable even with outlets shared, and across held-out outlets it
drops below the baseline entirely, so what a lexical model learns here is
outlet identity, not ideology.}
\label{tab:probe}
\end{table}

These are the lexical floor. The same probes on frozen encoder representations,
before and after our objective, are the core measurement of the project: the
question is whether the triplet term lowers publisher recoverability from the
representation \emph{without} lowering ROUGE.

\section{Background and Related Work}

\paragraph{NeuS} \citep{lee2022neus} introduces the task, the AllSides-derived
corpus of left/centre/right article triples with human-written neutral
``roundup'' summaries, and the framing-bias metric. It fine-tunes BART
\citep{lewis2020bart} on the concatenated sources and reports that a multi-task
auxiliary objective over the encoder helps, but the auxiliary task is
classification, and nothing ties the representation to event identity.

\paragraph{We Can Detect Your Bias} \citep{baly2020detect} predicts
article-level political ideology and uses a triplet objective to strip publisher
identity from a representation while preserving ideology. It also contributes the
media-based split protocol, in which whole outlets are held out. We invert the
objective: preserve \emph{event} identity while removing \emph{source}
information.

\paragraph{Multi-document event relation graphs} \citep{lei2025multidoc} target
the same framing-bias mitigation goal on the same corpus by building an explicit
event relation graph and reasoning over it with an LLM. Our approach is the
representation-level counterpart: rather than externalising event structure into
a graph, we constrain the encoder's latent space to be organised by event. We
regard this as the closest prior work and will compare against it directly.

\paragraph{I-JEPA} \citep{assran2023ijepa} predicts one view's latent
representation from another rather than reconstructing raw content; this
motivates our label-efficient pretraining extension, where the ``views'' are two
outlets' reports of the same event.

\paragraph{DiverseSumm} \citep{huang2024diversesumm} proposes QA-based
evaluation of whether a multi-document summary retains information from its
sources.

\section{Contributions}

\begin{enumerate}\itemsep2pt
\item \textbf{An event-aware triplet objective} for multi-source neutral
summarization, in which the positive shares an event with the anchor and the
negative shares a \emph{publisher}. The shared publisher is the mechanism: it is
the only configuration in which satisfying the margin requires the encoder to
stop relying on outlet style. This addresses the first shortcoming: that
standard training leaves the encoder's organising principle unconstrained.

\item \textbf{A length-normalised framing-bias metric}, together with the
evidence that the standard metric is a length proxy. We will report both, and
show what conclusions change. This addresses the third shortcoming, and makes
neutrality-by-deletion visible rather than rewarded.

\item \textbf{A publisher-disjoint evaluation protocol} for this corpus: a
stance-stratified split in which no outlet appears on two sides, plus a leakage
probe measured on both the standard and the disjoint split. This addresses the
fourth shortcoming: that the released split cannot ask the generalisation
question.

\item \textbf{A reproducible data pipeline.} The released NeuS split files have
their publisher and stance labels stripped. We recover them at $100\%$ coverage
by normalising headlines against the union of all three crawl files; the
obvious single-file join silently misses $13\%$ of examples. We will release
this alongside our code.
\end{enumerate}

\section{Proposed Methodology}

\paragraph{Backbone.} A pretrained sequence-to-sequence model (BART-large). The
$k$ articles for an event are concatenated into a single source sequence and
decoded into a neutral summary under the usual token-level cross-entropy loss
against the human-written roundup.

\paragraph{Event-aware triplet objective.} For an anchor article $a$, we sample a
positive $p$ (a different outlet's report of the same event) and a negative
$n$ (the same outlet's report of a different event). Each is passed through the
shared encoder and pooled at the EOS position to give $f(\cdot)$. The training
signal is
\begin{align*}
\mathcal{L} = \mathcal{L}_{\text{gen}} + \lambda \cdot \max\big(&0,\;
  \|f(a)-f(p)\| \\
  &- \|f(a)-f(n)\| + m\big),
\end{align*}
with margin $m$ and weight $\lambda$ swept on the validation set. Because $a$ and
$n$ always share a publisher, the only way to satisfy the margin consistently is
to stop encoding publisher style and organise around the event instead.

Structurally this attaches at the same point as NeuS's own auxiliary-task hook,
which adds a weighted auxiliary loss over a pooled encoder state; we replace
their binary classification head with the triplet term and reimplement on a
current \texttt{transformers} stack rather than their pinned 2021 fork.

\paragraph{Sampling and training details.} Triples are mined offline: for each of
the 9{,}282 articles we draw two positives and two negatives, discarding anchors
whose outlet appears in only one story, which yields 18{,}368 triples. The
triplet term is fed as a second stream alongside the generation batch, mirroring
the reference implementation's separate auxiliary-input path, so the two losses
need not share a batch size. We start from the hyperparameters NeuS reports
(learning rate $3\!\times\!10^{-5}$, effective batch 16, source length 512, target
length 250) and sweep only $\lambda$ and $m$. Training is short: 2{,}452
examples over five epochs is roughly 3{,}000 optimiser steps, under an hour of
single-GPU time per configuration, which is what makes the ablation grid and the
multi-seed reporting affordable.

\paragraph{Extension: cross-source pretraining.} Following the JEPA principle, we
pretrain by predicting one outlet's representation of an event from another
outlet's representation of the same event, using unlabelled same-event pairs from
WCEP-10, which needs neither publisher labels nor neutral targets, and so
suits this objective exactly. The question is whether the decoder can then be
fine-tuned on a fraction of the labelled summaries and still match a fully
supervised model. This extension is contingent on compute and will be dropped
before the core result if necessary.

\section{Dataset}

Our primary corpus is the NeuS release \citep{lee2022neus}: 3{,}095 AllSides
stories, each pairing three articles (one left, one centre, one right) with a
human-written neutral roundup, spanning 2016--2021 and 250 distinct outlets. It
is the only corpus available to us that carries an event grouping, a publisher
label and a neutral target on the same example. From it we derive 9{,}282
publisher-labelled articles and 18{,}368 training triples; 153 of the 250 outlets
appear in at least two distinct stories, which makes 9{,}188 of 9{,}285 article
slots usable as a same-publisher negative.

Two properties we flag honestly. First, AllSides publishes an \emph{excerpt} per
outlet rather than the full article (a hard cap of 95 words, median 85), so
a three-source input is $\approx$241 words and fits comfortably within BART's
1{,}024-token window, but there is correspondingly little event content to
retain. Second, at 2{,}452 training examples the corpus is small for fine-tuning
BART-large; we will report mean and spread over at least three seeds rather than
single runs.

We use two further corpora in supporting roles. WCEP-10
\citep{ghalandari2020wcep} provides 10{,}200 genuine multi-document event
clusters of $\approx$9 articles each, far better event structure and scale
than NeuS, but its preprocessing stripped the publisher labels and its 25-word
targets were never written to be neutral, so we use it only for the pretraining
extension. AllSides / Article-Bias-Prediction \citep{baly2020detect} supplies
37{,}554 publisher-labelled articles for additional negatives and for its
prebuilt media-based split. We note that its bias labels are outlet-level: only 7
of its 491 sources ever carry more than one label, so ideology prediction and
publisher prediction are near-identical tasks there, and we will not report them
as independent axes.

\section{Evaluation Plan}

\paragraph{Event retention: question--answer coverage.} We generate
question--answer pairs from the source articles and measure how many remain
answerable from the generated summary, following the QA-based faithfulness
protocol of FEQA \citep{durmus2020feqa}. There are two ways to run this, and they
trade off against each other. (i) On our own NeuS-derived evaluation set: this
keeps the stance labels, and so is the only option that can ask whether coverage
is \emph{balanced} across left, centre and right sources; that is, whether the
model systematically drops facts that appear only in the right-leaning article.
It is the more expensive option in time and compute. (ii) On DiverseSumm
\citep{huang2024diversesumm}: purpose-built for this measurement, but its
articles carry no ideology labels, so it cannot detect that failure mode, and
running different evaluation axes on different corpora would make the comparison
unfair. We therefore plan on (i) and treat DiverseSumm as a secondary sanity
check.

\paragraph{Tonality.} We score generations with the NRC-VAD
\citep{mohammad2018vad} valence--arousal--dominance lexicon, reporting both the
framing-bias score as NeuS defines it and the length-normalised affect rate
introduced above, so that a score lowered by generating less is distinguishable
from one lowered by generating more neutrally.

\paragraph{Publisher leakage.} We train a publisher classifier on frozen
representations and on output text, and run it over both the source articles and
the generated summaries. We expect macro-F1 to fall, ideally to the random
baseline. This is measured on both the standard split and our publisher-disjoint
split.

\paragraph{Summary quality and ablations.} We train two models (a control BART
fine-tuned with the generation loss alone, and the same model with the triplet
term added) and compare every metric above plus ROUGE across the two, so that
any significant effect can be attributed to the objective rather than to the
backbone or the data. The $\lambda$ and margin sweeps give the intermediate
points between them.

\section{Feasibility and Risks}

The compute is modest: the full experimental grid (baseline, triplet sweep,
ablations, and reruns on the disjoint split) is roughly 20--25 single-GPU
hours at 16\,GB, which our cluster allocation covers comfortably. The data
pipeline, the metric, and the probes are CPU-only and already run.

Three risks we are managing explicitly. \textbf{Small data:} 2{,}452 training
examples will produce noticeable seed variance, so every reported number will be
a mean over at least three seeds with spread. \textbf{Weak event content:} the
85-word source excerpts limit what an event-retention evaluation can claim, so we
have demoted the QA check to a secondary measure rather than build the argument
on it. \textbf{A null result is possible:} the triplet term may reduce publisher
recoverability without improving summary neutrality, or may trade ROUGE for it.
We regard that as a reportable outcome; combined with the metric critique it
still constitutes a complete finding about what representation-level constraints
can and cannot buy in this task.

\section{Future Work}

Three directions fall outside the twelve weeks but follow directly from the
design.

\paragraph{Shared/unique factorisation.} The triplet term treats publisher
information as something to remove. A FactorCL-style decomposition
\citep{liang2023factorcl} would instead split the representation into a
component shared across all reports of an event and a component unique to each
outlet, keeping both and generating only from the shared part. That is a
strictly more informative object than a single scrubbed vector: the unique
component is itself a measurement of how a given outlet frames a given story.
We treat this as the natural successor to the present objective rather than a
competing one.

\paragraph{Full articles rather than excerpts.} Our sharpest limitation is that
AllSides publishes 85-word excerpts, which caps how much event content there is
to retain and therefore how much an event-retention metric can say. Recrawling
the source URLs for full text would remove that ceiling and let the QA
evaluation carry real weight, at the cost of a crawl and the licensing questions
that come with it.

\paragraph{Beyond BART and beyond US media.} BART-large is the right backbone
for a controlled study on 2{,}452 examples, but the objective is architecture
agnostic and could attach to any encoder. Separately, AllSides' left/centre/right
axis is a US-specific, outlet-level construct; whether an event-aware objective
transfers to media ecosystems with different partisan structure is an open
question our data cannot answer.

\section{Broader Impact}

\paragraph{Who defines neutral.} Our neutral targets are AllSides' human-written
roundups, and our stance labels are outlet-level judgements from the same
organisation. A model trained this way does not learn neutrality in the
abstract; it learns one editorial organisation's operationalisation of it. We
will state this framing explicitly wherever we report results rather than let
``neutral'' stand unqualified.

\paragraph{Neutrality by deletion is a harm, not just a metric artefact.} The
failure mode our metric critique identifies is also the ethically serious one. A
summariser that scores well by dropping whatever the sources disagree about
presents contested reporting as settled and removes exactly the information a
reader needs to judge the story. This is why we report the length-normalised
rate alongside the raw score, and why we regard a system that lowers framing
bias purely by saying less as having failed rather than succeeded.

\paragraph{Dual use of the leakage probe.} The probe we use to show that
publisher identity has been removed is, run the other way, a stylometric
attribution tool for news text. We are building it on published outlets with
public bylines, where attribution is not sensitive, but the same method applied
to anonymous or pseudonymous writing is a deanonymisation technique. We note
this because our contribution is partly the probe protocol itself.

\paragraph{Automation bias.} A summary presented as neutral invites more trust
than the same text presented as one outlet's account. Any deployment would need
to surface its sources rather than replace them.

\section{Conclusion}

We propose to constrain a multi-source summarisation encoder to organise its
latent space around events rather than publishers, using a triplet objective
whose negative deliberately shares a publisher with the anchor. The premise is
measured rather than assumed: on the NeuS corpus, publisher identity is
lexically recoverable at roughly twice its baseline while stance does not
transfer across held-out outlets at all. Our audit also found that the field's
framing-bias metric is largely a length proxy and that the standard split shares
almost all of its test outlets with training, so we contribute a
length-normalised metric and a publisher-disjoint protocol alongside the
objective. Whether the triplet term buys lower publisher recoverability without
costing ROUGE is the question the project answers; a null result, read against
the metric critique, is still a complete finding about what representation-level
constraints can and cannot do for this task.

\section{Timeline}

\begin{table}[h]
\centering\small
\begin{tabular}{@{}ll@{}}
\toprule
\textbf{Weeks} & \textbf{Milestone} \\
\midrule
1--2   & Data pipeline, splits, probe protocol \\
       & \emph{(complete)} \\
3--4   & NeuS baseline reproduced; metric \\
       & analysis written up \\
5--6   & Triplet objective implemented; \\
       & $\lambda$/margin sweep started \\
\midrule
\multicolumn{2}{@{}l@{}}{\emph{Mid-submission: baseline + first triplet}} \\
\multicolumn{2}{@{}l@{}}{\emph{result + metric analysis}} \\
\midrule
7      & Leakage probes on both splits \\
8--9   & Ablations, seed variance, QA check \\
10--11 & Cross-source pretraining extension \\
       & (compute permitting) \\
12     & Final report and code release \\
\bottomrule
\end{tabular}
\caption{Planned schedule. Weeks 1--2 are already complete. The
mid-submission checkpoint falls after week 6, by which point the baseline and a
first triplet result exist.}
\label{tab:timeline}
\end{table}

\bibliography{references}

\end{document}
